\documentclass[12pt,a4paper]{article}

\usepackage[utf8]{inputenc}
\usepackage[T1]{fontenc}
\usepackage[french]{babel}
\usepackage[a4paper,margin=2.2cm]{geometry}
\usepackage{graphicx}
\usepackage{float}
\usepackage{xcolor}
\usepackage{listings}
\usepackage{hyperref}

\hypersetup{
    colorlinks=true,
    linkcolor=blue,
    urlcolor=blue
}

\definecolor{codegray}{gray}{0.95}
\definecolor{commentgreen}{rgb}{0.0,0.45,0.0}

\lstset{
    language=JavaScript,
    backgroundcolor=\color{codegray},
    basicstyle=\ttfamily\small,
    keywordstyle=\color{blue},
    commentstyle=\color{commentgreen},
    stringstyle=\color{red!70!black},
    numbers=left,
    numberstyle=\tiny\color{gray},
    numbersep=8pt,
    frame=single,
    breaklines=true,
    showstringspaces=false,
    tabsize=2
}

\title{Compte rendu -- Structuration de l'API Express}
\author{Nom et prénom : \rule{7cm}{0.4pt}}
\date{10 octobre 2026}

\begin{document}

\maketitle
\tableofcontents
\newpage

\section{Objectifs}

Le travail réalisé consiste à structurer une API Express selon le principe de séparation des
responsabilités (Separation of Concerns, SoC). La logique des domaines \texttt{articles} et
\texttt{users} est placée dans des contrôleurs dédiés, tandis que les routeurs associent les
méthodes HTTP aux fonctions correspondantes.

Les fonctionnalités réalisées sont les suivantes :
\begin{itemize}
    \item extraction de la logique des articles dans un contrôleur et un routeur ;
    \item ajout des opérations PUT et DELETE pour les articles ;
    \item extraction de la logique des utilisateurs dans un contrôleur et un routeur ;
    \item validation réutilisable du nom et de l'adresse email dans un module \texttt{utils}.
\end{itemize}

\section{Réponses aux encadrés « Prédisez »}

\subsection{Prédiction 1 -- La réponse des articles après la refactorisation}

\textbf{Prédiction :} la réponse de \texttt{GET /api/articles} ne changera pas après la
refactorisation.

\textbf{Justification :} le déplacement de la logique depuis \texttt{server.js} vers
\texttt{controllers/articleController.js} ne modifie ni les données initiales, ni le format de
la réponse, ni l'URL publique. Le routeur est simplement monté avec le même préfixe
\texttt{/api/articles}. Le client extérieur ne voit donc aucune différence ; seule
l'organisation interne du code change.

\subsection{Prédiction 2 -- Type du paramètre d'identifiant}

\textbf{Prédiction :} dans une route comme \texttt{/api/articles/2} ou \texttt{/api/users/2},
\texttt{req.params.id} est une chaîne de caractères, donc sa valeur est \texttt{"2"} et non le
nombre \texttt{2}.

\textbf{Justification :} les paramètres de route Express sont reçus sous forme de chaînes. Le
code utilise donc \texttt{Number(req.params.id)} avant de comparer l'identifiant avec les
identifiants numériques stockés dans les tableaux. Cette conversion permet à \texttt{find} ou
\texttt{findIndex} de retrouver correctement la ressource.

\section{Captures Postman -- Domaine articles}

Les figures suivantes sont des emplacements réservés. Il faut remplacer le contenu vide de
\texttt{\\includegraphics} par le chemin de chaque capture d'écran avant la compilation finale.

\subsection{Mise à jour réussie}

Requête : \texttt{PUT /api/articles/2}. Le serveur doit répondre avec le statut \texttt{200 OK}
et retourner l'article modifié.

\begin{figure}[H]
    \centering
    \includegraphics[width=1\textwidth]{}
    \caption{Mise à jour réussie de l'article 2 avec PUT /api/articles/2.}
    \label{fig:article-put-success}
\end{figure}

\subsection{Recherche d'un article inexistant}

Requête : \texttt{GET /api/articles/99}. Le serveur doit répondre avec le statut
\texttt{404 Not Found}.

\begin{figure}[H]
    \centering
    \includegraphics[width=1\textwidth]{}
    \caption{Erreur 404 lors de la recherche d'un article inexistant.}
    \label{fig:article-get-404}
\end{figure}

\subsection{Suppression réussie}

Requête : \texttt{DELETE /api/articles/1}. Le serveur doit répondre avec le statut
\texttt{200 OK} et confirmer la suppression.

\begin{figure}[H]
    \centering
    \includegraphics[width=1\textwidth]{}
    \caption{Suppression réussie de l'article 1 avec DELETE /api/articles/1.}
    \label{fig:article-delete-success}
\end{figure}

\section{Captures Postman -- Domaine users}

\subsection{Liste complète des utilisateurs}

Requête : \texttt{GET /api/users}. La réponse contient les deux utilisateurs fictifs avec leurs
champs \texttt{id}, \texttt{name}, \texttt{email} et \texttt{role}.

\begin{figure}[H]
    \centering
    \includegraphics[width=1\textwidth]{}
    \caption{Liste complète des utilisateurs avec GET /api/users.}
    \label{fig:user-get-all}
\end{figure}

\subsection{Création valide}

Requête : \texttt{POST /api/users} avec un corps JSON valide :

\begin{lstlisting}[language=json]
{
  "name": "Alex",
  "email": "alex@example.com",
  "role": "reader"
}
\end{lstlisting}

Le serveur doit répondre avec le statut \texttt{201 Created}.

\begin{figure}[H]
    \centering
    \includegraphics[width=1\textwidth]{}
    \caption{Création réussie d'un utilisateur avec POST /api/users et statut 201.}
    \label{fig:user-post-success}
\end{figure}

\subsection{Email invalide}

Requête : \texttt{POST /api/users} avec un email ne contenant pas les caractères requis, par
exemple \texttt{ayaexample.com}. Le serveur doit répondre avec le statut \texttt{400 Bad
Request} et un message explicite.

\begin{figure}[H]
    \centering
    \includegraphics[width=1\textwidth]{}
    \caption{Erreur 400 lors de la création d'un utilisateur avec un email invalide.}
    \label{fig:user-post-invalid-email}
\end{figure}

\section{Architecture de l'application}

Le fichier \texttt{server.js} monte les deux routeurs :

\begin{lstlisting}
app.use('/api/articles', articleRoutes);
app.use('/api/users', userRoutes);
\end{lstlisting}

Le routeur des articles utilise le contrôleur des articles et le routeur des utilisateurs utilise
le contrôleur des utilisateurs. Les données sont conservées en mémoire ; elles sont réinitialisées
au redémarrage du serveur.

\section{Code source demandé}

\subsection{controllers/userController.js}

\lstinputlisting{controllers/userController.js}

\subsection{routes/userRoutes.js}

\lstinputlisting{routes/userRoutes.js}

\subsection{utils/validators.js}

\lstinputlisting{utils/validators.js}

\subsection{server.js}

\lstinputlisting{server.js}

\section{Conclusion}

La logique des articles et des utilisateurs est maintenant organisée par domaine. Les contrôleurs
contiennent la logique métier, les routeurs assurent l'aiguillage des requêtes et \texttt{server.js}
conserve la configuration globale de l'application. Le module \texttt{validators.js} évite de
répéter les règles de validation et peut être réutilisé par d'autres contrôleurs.

\end{document}